\hspace{3mm}

\centerline{\bf \Huge  Формула Пика. Клетчатая геометрия.}
\centerline{\bf \Huge Уровень: Hard.}

\begin{flushright}
    \scriptsize{}
\end{flushright}


\problem{0}
Нарисуйте треугольник площади $\dfrac{1}{2}$, у которого все стороны больше 5, а вершины лежат в узлах сетки.

\problem{1} 
Докажите, что найдётся такая прямая на координатной плоскости, проходяшая через два узла и найдётся к ней такой узел, не лежащий на этой прямой, что расстояние между этим узлом и прямой будет меньше $\dfrac{1}{2024}$.

\problem{2} 
Докажите, что при любом расположении квадрата $n\times n$ на плоскости он покроет не больше чем $(n+1)^2$ точек.

\problem{3}
Точку M внутри треугольника соединили с его вершинами, в результате треугольник разбился на три равновеликие части. Докажите, что M — точка пересечения медиан треугольника.

\textit{Определение.} 
Треугольник с вершинами в целых точках назовём \textit{элементарным}, если внутри него и на его границе нет целых точек, кроме вершин.

\problem{4} 
Докажите, что элементарный треугольник не может быть остроугольным.

\problem{5} 
Докажите, что если три вершины параллелограмма лежат в целых точках, то и четвёртая вершина расположена в целой точке.

\problem{6} 
Докажите, что если площадь параллелограмма равна 1 и его вершины лежат в узлах, то любой из диагоналей он разбивается на два элементарных треугольника. 

\problem{7} 
Для треугольника на плоскости назовём скачком следующую операцию: одна из его вершин заменяется на симметричную ей относительно другой вершины.

(а) Докажите,что элементарный треугольник будет переходить в элементарный.

(б) \textbf{Самостоятельно.} Докажите,что в конце концов получится прямоугольный треугольник.


\problem{8}
$\bigl[\textbf{I Олимпиада ЭлМШ, 8 класс, последняя задача}\bigr]$
Баир выписал на доску в порядке возрастания все рациональные дроби между нулём и единицей, знаменатели, которых не превосходят $n$. Пусть $\dfrac{a}{b}$ и $\dfrac{c}{d}$ - две соседние (несократимые) дроби, выписанные на доску. Докажите, что $|ad - bc| = 1$.

\problem{9}
$\bigl[\textbf{ТурГор 98/99, 10-11 класс, сложный тур, последняя задача}\bigr]$
Ладья, шагая по одной клетке, за 64 хода обошла все клетки шахматной доски и вернулась на исходную клетку. Докажите, что число ходов по вертикали не равно числу ходов по горизонтали.
